\item Está trabajando en un empleo de verano en una empresa que diseña sistemas de energía no tradicionales. La compañía trabaja en una planta de energía eléctrica propuesta que usaría el gradiente de temperatura en el océano. El sistema incluye una máquina térmica que funcionaría entre $20.0\text{ }^\circ\text{C}$ (temperatura del agua superficial) y $5.00\text{ }^\circ\text{C}$ (temperatura del agua a una profundidad de aproximadamente 1 km). 

(a) Su supervisor le pide que determine la efectividad máxima de dicho sistema. 
(b) Además, si la salida de energía eléctrica de la planta es $75.0\text{ MW}$ y opera a la máxima eficiencia posible teórica, debe determinar la velocidad de la energía que se toma desde el depósito caliente. 
(c) A partir de esta información, si una factura de electricidad de un hogar típico muestra un uso de $950\text{ kWh}$ por mes, su supervisor quiere saber cuántos hogares pueden tener la potencia de este sistema de energía que funciona a su máxima eficiencia. 
(d) A medida que la energía se extrae del agua tibia superficial para operar el motor, esta es remplazada por la energía absorbida por la luz solar en la superficie. Si la intensidad promedio absorbida por la luz del sol es de $650\text{ W/m}^2$ durante 12 horas de luz diurna en un día despejado, debe encontrar el área de la superficie del océano que es necesaria para que la luz solar remplace la energía absorbida por el motor. 
(e) A partir de esta información, debe determinar si hay suficiente superficie oceánica en la Tierra para usar dichos motores y satisfacer las necesidades eléctricas de todos los hogares asociados con la población de la Tierra. Suponga que la energía utilizada para un hogar en el inciso (c) es un promedio en todo el planeta. 
(f) En vista de los resultados en este problema, su supervisor ha pedido su conclusión en cuanto a si vale la pena desarrollar tal sistema. Tenga en cuenta que el “combustible” (luz solar) es gratis.
